\section{从数据到矩阵的实现细节}
\subsection{局部坐标尺度与链式法则}
设子域矩形中心为$x_i^c$，两个方向边长为$h_{i,x},h_{i,y}$，局部坐标为$\xi=(x-x_i^c)/h_{i,x}$和$\eta=(y-y_i^c)/h_{i,y}$。对$z=w_x\xi+w_y\eta+b$，特征的物理导数为
\[
\partial_x\tanh z=\frac{w_x}{h_{i,x}}(1-\tanh^2z),\quad
\partial_{xx}\tanh z=-2\frac{w_x^2}{h_{i,x}^2}\tanh z(1-\tanh^2z).
\]
另一方向同理。这里的中心化坐标与前文单位方形坐标减去$1/2$等价；直线子域取$h_{i,x}=.5,h_{i,y}=1$。若遗漏边长倒数，内部算子和通量均会改变；即使一个对称算例偶然给出较小误差，算法也不再求解原方程。本文把链式因子放在特征模块统一处理，界面组装不再重复缩放。

局部坐标映射与权重范围共同确定物理变化尺度。把同一个随机参数范围用于大小不同的子域，会产生不同的物理频率；因此报告scale时还必须说明是否映射到参考矩形。花形近似复现采用原物理坐标，不能直接把其参数1.6与直线局部映射中的scale2视为相近设置。两者即使都用tanh，矩阵列和导数的幅值也可能不同。

激活函数接近饱和时，一阶和二阶导数趋小，可能降低某些PDE观测的分辨能力。权重过小则多个特征在有限区域上接近低阶展开，容易近相关。这些都是范围敏感性的机制解释，本文没有据此推导某个普适最优范围；具体选择仍由事先登记的模型和扫描结果说明。

\subsection{装配维度、数据顺序及符号核对}
输出系数按两个子域拼接为$c=(c_1^T,c_2^T)^T$。左内部行只作用于$c_1$，右内部行只作用于$c_2$；边界按所在子域赋值。界面值行是$(-\Phi_1,\Phi_2)$，通量行为$(-\beta_1\partial_n\Phi_1,\beta_2\partial_n\Phi_2)$。右端分别放置$g_u$和$g_q$。若更换跳跃定义，整行和右端同时变号不改变约束；只改其中一侧则会求解另一个问题。

内部行、边界行及界面行分别保留组信息后再叠置，这允许单独计算每组残差并核对行数。所有组必须使用相同的系数列顺序；切分预测系数时也使用相同索引。仅检查最终矩阵可乘并不足以发现子域列交换，因为这种错误往往仍有合法维度。

制造解检验包含两层。第一层在任意非配置点计算解析PDE和跳跃，检查模型数据自身是否一致。第二层将特征导数与独立差分近似比较，检查链式因子和符号。模型一致性正确不意味着特征模块正确，反之亦然，因此两类测试分开。通量不恒零及非零跳跃的构造，可使遗漏扩散系数或漏写一侧约束更容易暴露。

\section{QR、SVD与正则化的不同作用}
\subsection{满列秩时的QR解释}
若$B$满列秩，薄QR分解为$B=QR$，其中$Q^TQ=I$且$R$可逆。将数据分解为$QQ^Td+(I-QQ^T)d$，两部分正交，得到
\[
\norm{Bc-d}^2=\norm{Rc-Q^Td}^2+\norm{(I-QQ^T)d}^2.
\]
第二项与$c$无关，因此最优系数由上三角方程$Rc=Q^Td$给出。这一推导说明QR避免形成$B^TB$；但当$R$本身接近奇异时，QR也不能凭空消除数据敏感性。实际秩亏问题可使用带列选主元的QR或SVD，本文采用SVD型实现以便统一报告奇异值和截断。

SVD把相关方向明确分解，允许以阈值舍去难以可靠区分的成分。然而阈值舍弃的可能是噪声放大方向，也可能包含逼近真解所需的信息。因此数值秩是求解约定与矩阵谱共同确定的结果，不能只按一次阈值所得秩给出特征空间的绝对数学维数。

\subsection{岭的偏差与残差权衡}
对$\alpha>0$，岭解满足
\[
c_\alpha=\sum_j\frac{\sigma_j}{\sigma_j^2+\alpha}(q_j^Td)v_j.
\]
与不截断伪逆相比，第$j$个方向乘以$\sigma_j^2/(\sigma_j^2+\alpha)$，小奇异方向收缩更强。若数据恰为$Bc_*$，则在可辨识方向上仍存在收缩偏差；如果$c_*$在这些方向上具有重要分量，训练稳定性改善也可能伴随更大的真解误差。实验中强岭导致误差增加，与这种偏差机制相容，但并不能反推出全部误差都由岭引起。

代码通过增广矩阵$[B;\sqrt\alpha I]$求解岭目标。增广矩阵的秩与原数据矩阵的有效秩不同；前者因正则化可变为满列秩，后者仍说明原观测可以分辨多少方向。论文所列data rank针对原数据矩阵，避免把正则化增加的对角行误认为采样信息增加。

\subsection{有限精度与实际求解的边界}
双精度机器精度约为$2.22\times10^{-16}$。当矩阵条件数接近其倒数时，最坏方向的相对扰动可能非常敏感；但这个判断是风险尺度，不能直接预测每个右端的实际误差。数据是否激发最坏方向、截断是否舍弃该方向，以及预测矩阵是否放大对应系数，都会影响最终函数变化。

评价算法时需要区分组装误差、线性求解误差和函数求值误差。解析特征导数减少了差分截断，但仍有浮点运算；SVD型最小二乘避免正规方程，不保证高度相关列下的所有系数准确；大系数相消又可能使函数值的实际可辨位数减少。本文固定矩阵右端实验只研究其中一个环节，未完成对全部误差来源的统一量化。

\section{有限差分特征导数的误差机制}
对足够光滑的一维函数，中心二阶差分的截断误差为$h^2\phi^{(4)}(\xi)/12$。在简单浮点误差模型中，若函数值含有量级$\epsilon_{\rm mach}|\phi|$的扰动，除以$h^2$后会出现量级$\epsilon_{\rm mach}|\phi|/h^2$的放大。一个常见的解释性模型是
\[
E(h)\lesssim C_1h^2+C_2\epsilon_{\rm mach}h^{-2}.
\]
这里常数依赖特征、导数范围和求值过程，符号$\lesssim$表示机制尺度，并非本项目全部特征上已经认证的统一上界。两项平衡提示最优尺度依赖浮点精度及函数高阶导数，而非步长越小越好。

对tanh特征，高阶导数还包含权重的相应幂，因此不同参数范围可改变差分误差。进一步把差分误差放入最小二乘矩阵后，其对解的作用又受奇异值和截断影响。仅比较单个导数近似误差不足以完全解释PDE解误差；本文同时保留解析导数对照和步长扫描，正是为了避免只观察某一环节。

差分邻点可以跨越几何界面，但对局部网络进行差分时，求值对象始终是同一侧定义的光滑特征延拓，不能在邻点位置改用另一侧网络。否则差分混合了两侧物理不连续性，不再是该侧特征导数近似。花形实现使用各自网络的物理坐标特征，内部组只在对应区域布点，界面仍使用独立两侧迹值。
